%%%%%%%%%%%%%%%%%%%%%%%%%%%%%%%%%%%%%%%%%
% Lachaise Assignment
% LaTeX Template
% Version 1.0 (26/6/2018)
%
% This template originates from:
% http://www.LaTeXTemplates.com
%
% Authors:
% Marion Lachaise & François Févotte
% Vel (vel@LaTeXTemplates.com)
%
% License:
% CC BY-NC-SA 3.0 (http://creativecommons.org/licenses/by-nc-sa/3.0/)
% 
%%%%%%%%%%%%%%%%%%%%%%%%%%%%%%%%%%%%%%%%%

%----------------------------------------------------------------------------------------
%	PACKAGES AND OTHER DOCUMENT CONFIGURATIONS
%----------------------------------------------------------------------------------------

\documentclass{article}

\input{structure.tex} % Include the file specifying the document structure and custom commands

\usepackage{tabularx}

%----------------------------------------------------------------------------------------
%	ASSIGNMENT INFORMATION
%----------------------------------------------------------------------------------------

\title{Metodologías-2025: Evaluación Integral \#4} % Title of the assignment

\author{Luis Fernando González Avila} % Author name and email address

\date{Consejo Latinoamericano de Ciencias Sociales --- August 17, 2025} % University, school and/or department name(s) and a date

%----------------------------------------------------------------------------------------

\begin{document}

\maketitle % Print the title

%----------------------------------------------------------------------------------------
%	INTRODUCTION
%----------------------------------------------------------------------------------------

\section*{Operacionalización del concepto de \textbf{bienestar}} % Unnumbered section

La integración de dimensiones materiales, subjetivas y sociales es la constitución que se da del \textbf{bienestar}; y es por eso que su operacionalización requiere, como plantean Lazarsfeld (1979) y González Blasco (1998), traducir nociones abstractas en dimensiones observables, indicadores concretos y escalas de medición. Asi pues, el bienestar no se limita a la posesión de bienes materiales, sino que además, incluye percepciones de satisfacción y redes de apoyo social, que complementa la idea de una “realidad multidimensional” defendida por Cea D’Ancona (2001).\\

\noindent La dimensión \textit{material} contempla los recursos económicos y las condiciones de vida básicas. Aquí se incluyen indicadores como el ingreso mensual, el acceso a una vivienda digna, la disponibilidad de servicios básicos (agua, luz, internet) y la estabilidad laboral. Estas variables se miden en escalas de intervalo o dicotómicas, lo que permite establecer comparaciones objetivas entre sujetos.\\

\noindent La segunda dimensión es la \textit{subjetiva}, vinculada a la percepción individual sobre la propia vida. Tal como lo retoma D'Ancona (2001), la satisfacción general con la vida, la autopercepción del estado de salud y los sentimientos de felicidad constituyen indicadores fundamentales y para captarlos, se utilizan escalas tipo \textit{Likert}, que permiten graduar las percepciones en que van desde un extremo a otro, es decir, en un continuo.\\

\noindent Finalmente, la dimensión \textit{social/relacional} subraya la importancia que tienen los vínculos interpersonales y comunitarios. En palabras de Baranger (2009), la vida social constituye un elemento clave en la construcción del bienestar, ya que amplía las oportunidades de apoyo y pertenencia. Por ello, se consideran indicadores como la frecuencia de contacto con familiares o amigos, la participación en organizaciones comunitarias y el apoyo percibido en situaciones de crisis. Estos pueden medirse con escalas de razón, dicotómicas y ordinales.\\

\noindent La tabla mostrada a continuación sintetiza esta propuesta de operacionalizaciónn:\\

\begin{center}
	\begin{tabular}{| m{2cm} | m{2cm}| m{6cm} | m{4cm} |} 
		\hline
		Concepto & Dimensiones & Indicadores & Escala de Medición \\ [0.5ex] 
		\hline\hline
		Bienestar & Material & Ingreso mensual, acceso a vivienda digna, disponibilidad de servicios básicos, estabilidad laboral & Intervalo, nominal, dicotómica, ordinal \\ 
		\hline
		 & Subjetiva & Satisfacción con la vida, autopercepción de la salud, sentimiento de felicidad & Likert \\
		\hline
		 & Social o relacional & Frecuencia de contacto, participación comunitaria/social, apoyo percibido & Razón, dicotómica, ordinal \\ [1ex]
		\hline
	\end{tabular}
\end{center}

\noindent En síntesis, operacionalizar el bienestar implica construir un puente entre el plano conceptual y el empírico (Lazarsfeld, 1979). Esta propuesta ofrece un marco coherente y flexible para medir un fenómeno complejo, articulando lo objetivo, lo subjetivo y lo social como dimensiones inseparables de la calidad de vida.\\

\subsection*{Set de 12 preguntas para incluir en un formulario de encuesta}

\textit{Economía}
\begin{enumerate}
	\item ¿Sus ingresos actuales le permiten cubrir sus necesidades básicas (alimentación, vivienda, transporte, servicios)?
	\begin{itemize}
		\item Nada
		\item Poco
		\item Suficiente
		\item Completamente
	\end{itemize}
	\item ¿Qué tan estable considera su situación laboral o fuente de ingresos?
	\begin{itemize}
		\item Muy inestable
		\item Algo inestable
		\item Estable
		\item Muy estable
	\end{itemize}
	\item ¿Cuenta con algún tipo de ahorro o reserva económica para imprevistos?
	\begin{itemize}
		\item Si, suficientes
		\item si, pero limitado
		\item No cuento con ahorros
	\end{itemize}
\end{enumerate}
\textit{Salud}
\begin{enumerate}
	\item ¿Cómo calificaría su estado general de salud?
	\begin{itemize}
		\item Muy malo
		\item Malo
		\item Regular
		\item Bueno
		\item Muy bueno
	\end{itemize}
	\item ¿Con qué frecuencia realiza actividades físicas (caminar, ejercicio, deporte)?
	\begin{itemize}
		\item Nunca
		\item Ocasionalmente
		\item Algunas veces por semana
		\item Del diario
	\end{itemize}
	\item ¿Cuenta con acceso regular a servicios médicos cuando los necesita?
	\begin{itemize}
		\item Siempre
		\item Casi siempre
		\item Ocasionalmente
		\item Nunca
	\end{itemize}
\end{enumerate}
\textit{Psicológica}
\begin{enumerate}
	\item En general, ¿qué tan satisfecho está con su vida actualmente?
	\begin{itemize}
		\item Muy insatisfecho
		\item Insatisfecho
		\item Satisfecho
		\item Muy satisfecho
	\end{itemize}
	\item En las últimas dos semanas, ¿con qué frecuencia se ha sentido desanimado o con poca esperanza?
	\begin{itemize}
		\item Nunca
		\item Pocas veces
		\item Varias veces
		\item Casi todos los días
	\end{itemize}
	\item ¿Qué tan optimista se siente sobre su futuro?
	\begin{itemize}
		\item Pesimista
		\item Poco optimista
		\item Optimista
		\item muy optimista
	\end{itemize}
\end{enumerate}
\textit{Social}
\begin{enumerate}
	\item ¿Con qué frecuencia se reúne con familiares o amigos cercanos?
	\begin{itemize}
		\item Nunca
		\item Una vez al mes o menos
		\item Varias veces al mes
		\item Semanalmente o más
	\end{itemize}
	\item ¿Qué tan satisfecho está con el apoyo que recibe de sus relaciones personales más cercanas?
	\begin{itemize}
		\item Muy insatisfecho
		\item Insatisfecho
		\item Satisfecho
		\item Muy satisfecho
	\end{itemize}
	\item ¿Participa en organizaciones, grupos comunitarios o actividades sociales?
	\begin{itemize}
		\item Nunca
		\item Ocasionalmente
		\item Con frecuencia
		\item Muy activamente
	\end{itemize}
\end{enumerate}


%----------------------------------------------------------------------------------------

\end{document}
